\documentclass[11pt,a4paper,twoside,openany]{book}
\ifdefined\SdsTestEnglish
  \usepackage[english,nologo,pdfa]{sdsthesis}
\else\ifdefined\SdsTestKuten
  \usepackage[nologo,punctuation=kuten]{sdsthesis}
\else
  \usepackage[nologo]{sdsthesis}
\fi\fi

\thesissetup{
  type={回帰テスト},
  title={表紙と参照の確認},
  author={山田 太郎},
  affiliation={ソーシャル・データサイエンス学部},
  date={2026年3月 提出}
}
% Japanese mode must compile without any English metadata.
\ifdefined\SdsTestEnglish
  \thesissetup{
    etitle={Regression test},
    eauthor={Taro Yamada},
    eaffiliation={School of Social Data Science},
    edate={March, 2026}
  }
\fi
\begin{document}
\frontmatter
\maketitle
\begin{abstract}
  ABSTRACT-START.
  \newcount\abstractparagraph
  \abstractparagraph=0
  \loop
  This paragraph checks that a long abstract can break across pages.
  日本語と英語を含む長い概要を、ページからはみ出さずに表示します。
  The next paragraph must remain visible and inside the page margins.\par
  \advance\abstractparagraph by 1
  \ifnum\abstractparagraph<35\repeat
  ABSTRACT-END.
\end{abstract}
\tableofcontents
\mainmatter
\chapter{本文}\label{chap:body}
\section{本文の節}\label{sec:body}
Punctuation: 読点、と句点。全角の，と．も片方に統一される。
Math: $\mathbf{A}\mathbf{x} = \mathrm{d}\mathit{v} + \symbfit{y}$.
Forward: \cref{chap:appendix}; \cref{sec:appendix};
\cref{subsec:appendix}; \cref{eq:appendix};
\cref{fig:appendix}; \cref{tab:appendix}.
\begin{equation}\label{eq:body}x=1\end{equation}
\begin{theorem}\label{thm:test}A theorem.\end{theorem}
\begin{lemma}\label{lem:test}A lemma.\end{lemma}
\begin{definition}\label{def:test}A definition.\end{definition}
\begin{proposition}\label{prop:test}A proposition.\end{proposition}
\begin{proof*}A proof.\end{proof*}
\cref{thm:test}; \cref{lem:test}; \cref{def:test}; \cref{prop:test}.
\begin{appendices}
  \chapter{追加の検証}\label{chap:appendix}
  \section{付録の節}\label{sec:appendix}
  \subsection{付録の項}\label{subsec:appendix}
  Backward: \cref{chap:body}; \cref{sec:body}; \cref{eq:body}.
  Self: \cref{chap:appendix}; \cref{sec:appendix}.
  \begin{equation}\label{eq:appendix}x=2\end{equation}
  \begin{equation}\label{eq:appendix-two}x=3\end{equation}
  \begin{equation}\label{eq:appendix-three}x=4\end{equation}
  Range: \crefrange{eq:appendix}{eq:appendix-three}.
  Multiple: \cref{eq:body,eq:appendix}.
  \begin{figure}[htbp]
    \centering\rule{30mm}{10mm}
    \caption{付録の図}\label{fig:appendix}
  \end{figure}
  \begin{table}[htbp]
    \caption{付録の表}\label{tab:appendix}
    \centering\begin{tabular}{cc}A&B\\1&2\end{tabular}
  \end{table}
  \chapter{次の付録}\label{chap:appendix-b}
  \begin{equation}\label{eq:appendix-b}x=5\end{equation}
\end{appendices}
\end{document}
